\section{Observables and the classification rule}
\label{app:observables}

\paragraph{Conventions.} We use the signed class relation
$s_{IJ}=\kappa_I\kappa_J\in\{-1,+1\}$. A $\{0,1\}$ indicator, which counts only
same-class pairs, cannot represent anti-correlation: under perfect reverse
alignment it returns $-(N/2-1)/(N-1)$ rather than $-1$, and its range depends on
$N$. Since half the diagram has the field favouring the out-group, the signed
indicator is the appropriate one. Directed trust has $t_{I|I}\equiv+1$, which the
closed form below assumes.

\paragraph{Closed form for the balance aggregates.} Enumerating triples costs
$O(N^3)$ per society, which is prohibitive across a replicated grid. For any $M$
with unit diagonal,
%
\begin{equation}
  \sum_{\substack{(I,J,K)\\ \text{distinct, ordered}}} M_{IJ}M_{JK}M_{KI}
  = \operatorname{tr}(M^{3}) - 3\Bigl(\textstyle\sum_{I,J}M_{IJ}M_{JI} - N\Bigr) - N ,
\end{equation}
%
and each unordered triple appears six times, three as the forward cycle and three
as the reverse. For symmetric $M$ this gives the ideological balance; for the
asymmetric trust matrix the two orientations are exactly the two terms of the
affective balance, so one expression serves both. The identity is verified
against enumeration in the test suite, for continuous and sign matrices alike.

\paragraph{Finite-size values of the sanity checks.} Two configurations are worth
recording because the natural expectation is wrong at finite $N$. If every pair
trusts every other and all opinions align, then with balanced classes
$\CcT=\CcO=-1/(N-1)$, not $0$ --- averaging $s_{IJ}$ over unordered pairs of a
balanced split leaves a deficit of one pair's worth. And the polarisation
\eqref{eq:polarisation} cannot reach $1$: zeroing the diagonal of a rank-one
$ss^{\!\top}$ leaves a spectrum $N-1$ once and $-1$ with multiplicity $N-2$ after
centring, so $P_{\max}(N)=(N-1)^2/\bigl((N-1)^2+N-2\bigr)$, which is $0.94$ at
$N=16$ and $0.995$ at $N=200$. Both values are asserted in the tests, and
$P_{\max}$ is divided out so that polarisation is comparable across $N$.

\paragraph{The classification rule.} With thresholds $\theta_c$ on $|\CcT|$,
$\theta_p$ on the standardised polarisations and $\theta_b$ on $\Bt$, each
replicate is assigned in this fixed order: no class correlation and no
polarisation in either sector $\to$ \emph{weakly structured}; no class
correlation but polarisation present $\to$ \emph{polarized, class-uncorrelated};
class correlation of the wrong sign $\to$ \emph{counter-aligned, frustrated};
class correlation with opinion following class $\to$ \emph{discriminatory,
ideological}; otherwise $\to$ \emph{discriminatory, class-dominant}. Thresholds
are three null standard deviations of the corresponding statistic measured at
$d=0$ at the same $N$, so they carry the finite-size noise floor and need no
re-tuning by eye. The modal label over replicates is reported together with the
agreement fraction, and regime areas are recomputed at $\{2/3,1,3/2\}$ times each
threshold.

We considered clustering the order-parameter vectors instead and rejected it:
cluster identity is not stable across $N$, across agenda size, or between a
baseline and its controls, and comparing exactly those things is the purpose of
the finite-size and control figures.
